\section{Integración de los pipelines}
\label{sec:celdas_vs_pixeles}

% Explicar la diferencia entre el procesamiento del nnELM (grilla de
% celdas) y el pipeline de EEG (fijaciones en píxeles).
% Justificar el uso de los scanpaths de Damián como referencia.
% Describir la estrategia de asignación de valores de celda a fijaciones.
% Tratar el caso de fijaciones consecutivas en la misma celda.
% Mencionar los criterios de exclusión aplicados.

La incorporación de las features extraídas del nnELM al modelo de TRFs por regresión regularizada requiere compatibilizar dos pipelines que difieren en dos aspectos: el algoritmo de detección de fijaciones y la resolución espacial.

\subsubsection{Algoritmos de detección de fijaciones}

El pipeline de EEG utiliza el algoritmo de Engbert y Mergenthaler (2006) para detectar fijaciones y sacadas a partir de la señal del eye-tracking. El modelo nnELM, en cambio, fue diseñado para generar sus propios scanpaths internos a partir de su mecanismo de decisión. Dado que el análisis del EEG requiere que los eventos estén alineados con la señal cerebral, se forzó al modelo nnELM a utilizar cómo input las coordenadas de fijación del pipeline del EEG, en lugar de generar su propio scanpath. De esta manera, cómo resultado cada fijación del scanpath de \citet{care2025} obtuvo los valores de las features calculadas por el nnELM.

\subsubsection{Resolución espacial: celdas vs. píxeles}

El modelo nnELM opera sobre una representación espacial discreta de la imagen, dividida en una grilla de $24 \times 32$ celdas de $32 \times 32$ píxeles cada una. Como consecuencia, las features extraídas del modelo toman un valor por celda, no por píxel. El pipeline del EEG, en cambio, registra la posición de cada fijación a nivel pixel. Para asignar el valor de cada feature a cada fijación, se mapeó la coordenada de fijación a la celda correspondiente de la grilla.

Una consecuencia de esta discretización es que dos fijaciones consecutivas que caen dentro de la misma celda podrían recibir el mismo valor en las features extraídas del modelo nnELM. Sin embargo, esto sucede en las features estáticas del modelo (que no dependen del estado interno del modelo en cada paso). En las features dinámicas no sucede: entre fijación y fijación el modelo actualiza su estado interno, haciendo que el valor de las features cambien aunque caigan dentro de la misma celda. 


\begin{table}[h]
\centering
\small
\begin{tabular}{lrr}
\toprule
Feature & \% idénticas & Tipo \\
\midrule
\texttt{saliencia\_media}              & 100.0\% & estática \\
\texttt{saliencia\_mediana}            & 100.0\% & estática \\
\texttt{saliencia\_pixel}              &  54.9\% & estática \\
\texttt{similitud\_media}              & 100.0\% & estática \\
\texttt{similitud\_mediana}            & 100.0\% & estática \\
\texttt{similitud\_pixel}              &  37.9\% & estática \\
\texttt{entropia\_prior\_saliencia}    & 100.0\% & estática \\
\midrule
\texttt{competencia\_densidad}         &  66.8\% & dinámica \\
\texttt{competencia\_densidad\_norm}   &  66.8\% & dinámica \\
\texttt{distancia\_grilla}             &  58.1\% & dinámica \\
\texttt{distancia\_pixels}             &  58.1\% & dinámica \\
\texttt{gap\_entropia}                 &   3.1\% & dinámica \\
\texttt{evidencia\_visual}             &   0.0\% & dinámica \\
\texttt{competencia\_fp\_var}          &   0.0\% & dinámica \\
\texttt{competencia\_fp\_media}        &   0.0\% & dinámica \\
\texttt{competencia\_fp\_max}          &   0.0\% & dinámica \\
\texttt{posterior}                     &   0.0\% & dinámica \\
\texttt{entropia\_posterior}           &   0.0\% & dinámica \\
\texttt{ganancia\_informacion}         &   0.0\% & dinámica \\
\texttt{competencia\_ratio}            &   0.0\% & dinámica \\
\bottomrule
\end{tabular}
\caption{Proporción de pares de fijaciones consecutivas en la misma celda 
($N = 2{,}908$) con valores idénticos para cada feature. Las features estáticas 
no dependen de la secuencia de fijaciones, por lo que presentan redundancia 
total o parcial. Las features dinámicas se actualizan en cada paso del modelo, 
resultando en valores distintos incluso cuando dos fijaciones caen en la misma celda.}
\label{tab:pares_celda}
\end{table}